\begin{definition}[The original global fraction cover]
\label{def:whole-original-global-fraction-open-cover}
\lean{EnvelopeProperty.BlowupHomogeneousTautologicalLine.specifiedOriginalGlobalFractionOpenCover}
\leanok
\uses{thm:resume-specified-generator-fraction-cover}
Let $X$ be integral, let $j$ and the homogeneous ideal $J$ be the original data, and let $F_0,\ldots,F_m$ generate $J$. The actual native glued blowup of $D=J\mathcal O_X$ has an actual open cover indexed by nonempty original coordinate-subordinate affine opens and their nonzero original pulled generator coefficients. Each covering map is literally the actual native fraction inclusion followed by the actual native affine-blowup chart inclusion. No affine-domain, nonzero-ideal, cover or open-map certificate is an input.
\end{definition}
\begin{proof}
A nonempty open of an integral scheme has a domain section ring, giving the needed original affine ring internally. On each subordinate affine open, the original pulled generators span the actual original ideal, so their nonzero-generator native fraction charts cover its whole actual affine blowup. Both the fraction inclusion and the original global chart inclusion are open immersions. For a point of the global blowup, choose an original coordinate open containing its image under $j\mu$, then a subordinate affine neighbourhood of its image in $X$. This affine is nonempty because it contains that image point. The native global chart image is the inverse image of the affine neighbourhood, so the original point lifts to that actual affine blowup. Its actual finite generator fraction cover supplies a native fraction preimage. Composing the two original maps gives the required preimage of the global point. This also proves the zero-ideal case: the native empty affine blowups have empty nonzero-generator covers, and the same pointwise argument gives the genuine empty global cover.
\end{proof}
